\documentclass[10pt,a4paper]{amsart}
\usepackage[margin=19mm]{geometry}
\usepackage{amsmath,amssymb,mathtools}
\usepackage[scr=rsfs,cal=euler]{mathalfa}
\usepackage{xfrac}
\usepackage[unicode,hidelinks,bookmarks=false]{hyperref}
\newcommand{\ie}{i.e.\ }
\newcommand{\cf}{cf.\ }
\newcommand{\resp}{respectively\ }
\newcommand{\Subsection}{\subsection}
\providecommand{\nobreakdash}{\nobreak\hbox{-}}
\setlength{\emergencystretch}{2em}
\multlinegap0pt
\usepackage{amssymb}
\usepackage{enumerate}
\newtheorem{lemma}{Lemma}
\title{Logical Undefinability of the Generalized Collatz Transition Relation in B\"uchi Arithmetic}
\author{Madhav Dhiman, Rohan Pandey}
\date{Source: arXiv:2601.12772v2 (CC BY 4.0)}
\begin{document}
\maketitle

\noindent\emph{Source and licence.} Logical Undefinability of the Generalized Collatz Transition Relation in B\"uchi Arithmetic. Version arXiv:2601.12772v2 (CC BY 4.0), distributed by arXiv under the Creative Commons Attribution 4.0 International licence, http://creativecommons.org/licenses/by/4.0/, as recorded in the arXiv licence metadata for this exact version and shown on the abstract page https://arxiv.org/abs/2601.12772v2. Full licence notice: \texttt{licenses/arxiv-2601.12772-CC-BY-4.0.txt}.\par\smallskip

\noindent\textit{Editorial scope.} Counted unit: the article's lemma lem:pq\_definable (source0.tex lines 75--77). The article's complete original proof of it is delivered with it (source0.tex lines 79--81, one \texttt{proof} environment of 3 lines). No mathematics is written, repaired or re-derived: the statement and the proof are the article's own lines, in the article's own notation, and no retained line is substituted or re-flowed.  No further source-local context is needed: every cross-reference of every retained line resolves inside the two delivered spans.  Nothing was omitted from any retained span, so the package carries no omission record.  No retained line contains a citation, so the wrapper carries no bibliography and no bibliography entry.  No retained line contains a figure environment, an image-inclusion command or drawing code, so no image is delivered and the package declares no asset dependency.  Editorial formatting only, in addition: an AMS \texttt{amsart} preamble that restates the article's own declarations and macros used by the retained lines, namely the environments \texttt{lemma} on separate counters, together with the standard packages those lines need.  The retained environments print in this document as Lemma 1 in order of appearance, whereas the article prints the same items with the numbering of its own full document.  The complete native source of the article as distributed in the arXiv e-print for this exact version is bundled unchanged as \texttt{sources/192977/source0.tex}.
\begin{lemma} \label{lem:pq_definable}
Let $q > 1$ be an integer. The set of powers of $q$, defined as $P_q = \{q^y : y \in \mathbb{N}\}$, is first-order definable in $BA_q$.
\end{lemma}

\begin{proof}
In the structure $BA_q = \langle \mathbb{N}, +, 0, 1, V_q \rangle$, the function $V_q(x)$ returns the largest power of $q$ dividing $x$. The set of powers of $q$ is captured by the first-order predicate $Pow_q(x) \iff x > 0 \land V_q(x) = x$. Therefore, $P_q$ is trivially definable in $BA_q$.
\end{proof}

\end{document}
